\subsection{Cas 2 : patrimoine hors CTO de 300 000 euros}
\label{sec:cas2_sans_retrait}

La figure \ref{fig:sim2} illustre la différence d'héritage net dans un contexte patrimonial moins élevé, correspondant à un patrimoine hors enveloppes de 300 000 euros.

\begin{figure}[h!]
    \centering
    \includegraphics[width=\textwidth]{cas_2.png}
    \caption{Différence relative d'héritage net pour un investissement de 100 000 euros avant 70 ans, avec un patrimoine hors enveloppes de 300 000 euros.}
    \label{fig:sim2}
\end{figure}

Dans ce cas, l'impact du barème progressif est plus limité, le patrimoine demeurant longtemps dans les tranches basses (20 à 30 \%). Les résultats mettent en évidence que :
\begin{itemize}
    \item l'avantage comparatif de l'assurance-vie est moins marqué que dans le scénario précédent, les zones où cette enveloppe domine sont encore plus restreintes et le gain relatif semble encore plus faible;
    \item le CTO reste généralement plus favorable pour des horizons supérieurs à 15 ans ou en présence de rendements significatifs;
    \item l'exonération des plus-values au décès joue un rôle déterminant dans la supériorité du CTO.
\end{itemize}
